\chapter{结论}\label{chap:conclusion}

% ─────────────────────────────────────────────────────────────────────────────
\section{工作总结}\label{sec:summary}

本文提出并实现了一个基于物理信息神经网络（PINN）与大语言模型（LLM）的统一期权定价框架，将 BSM、CEV、Heston 三大经典期权定价模型纳入同一体系。

在 PINN 部分，三大模型共用统一的复合损失函数结构（PDE 残差 + 终值条件 + 边界条件），通过自动微分精确计算 PDE 残差，无需网格剖分。BSM 和 CEV 模型采用门控网络架构，浅层分支捕捉线性特征，深层分支捕捉平值附近的非线性特征；Heston 模型采用更深的全连接网络，并结合 ICPINN 硬约束输出变换，确保终值条件在机器精度下精确满足。

在 LLM 部分，引入 DeepSeek API 作为智能路由层，通过结构化 JSON 输出实现可靠的参数提取和模型选择，支持中英文自然语言输入，使非专业用户无需了解模型细节即可完成期权定价。

实验结果表明，BSM-PINN 在解析解验证下达到 MAE=0.056 的精度；CEV-PINN 正确捕捉了杠杆效应；Heston-PINN 在二维随机波动率场景下实现了可用的定价精度。LLM 路由层在测试用例上实现了 100\% 的模型选择准确率。

% ─────────────────────────────────────────────────────────────────────────────
\section{局限性与未来工作}\label{sec:future}

本文工作存在以下局限性，可在未来工作中改进：

\textbf{精度有待提升}：CEV 和 Heston 模型的相对误差（约 11\% 和 24\%）高于工程应用要求。可通过增加训练轮数、引入自适应配点采样（如 RAR 方法）、或采用更先进的网络架构（如 Fourier 特征网络）来改善精度。

\textbf{仅支持欧式看涨期权}：当前框架仅针对欧式看涨期权，未来可扩展至欧式看跌期权（通过 put-call parity 转换）和美式期权（需处理自由边界问题）。

\textbf{参数固定}：当前 PINN 模型针对固定参数集训练，若参数变化（如不同的 $\sigma$ 或 $K$）需要重新训练。未来可探索参数化 PINN（将参数作为额外输入），实现一次训练、多参数推断。

\textbf{LLM 参数提取的鲁棒性}：当前系统对模糊或不完整的自然语言输入处理能力有限。未来可引入多轮对话机制，让 LLM 主动向用户询问缺失参数。

\textbf{实时定价场景}：当前系统的主要延迟来自 LLM API 调用（约 1-2 秒）。对于高频交易场景，可考虑将 LLM 替换为轻量级的规则解析器，仅保留 PINN 的快速推断能力。
